\begin{abstract}
Recurrent latent reasoners can solve harder problems by iterating for longer, but the extra compute
only helps if the model also knows when to stop. We study this termination problem in a setting where
the answer to ``is the computation finished?'' is known exactly. For local recurrent models, a
light-cone argument assigns every input $x$ a lower bound $\LB(x)$ on the first step at which its answer
can be certified; the bound grows with depth rather than length, and a local, size-invariant automaton
attains it at every length. We introduce DepthShift, a benchmark with exact per-instance depth and
splits that vary depth and length independently, and evaluate fifteen stopping rules on one backbone
family. Jointly trained pooled gates (ACT, PonderNet, TRM-style heads) leak length: on deeper addition
inputs ACT is 99.8\% accurate while halting before $\LB(x)$ through a global shortcut, yet only 35.9\%
once inputs are also longer. Convergence rules stop before the bound on most deep inputs and are still
correct, so premature stopping and premature error must be measured separately. On addition,
per-position confidence halts at exactly $\LB(x)$ on every split, and the local completion heads we
propose match it without improving on it; on segmented parity the same heads learn the training clock
instead. Termination detection beyond the training depth appears to be limited less by the readout than
by training data that never show an unfinished computation at a late step. Code, frozen splits and all
runs are released.
\end{abstract}
